\section{The Goodhart Theorem for Proxy-Instrumental Alignment}
\label{sec:goodhart}

This section establishes the paper's novel formal contribution and
delivers the central alignment-payoff of the quantification claim.
Theorem~\ref{thm:fungibility_collapse} of
Section~\ref{sec:fungibility_collapse} placed standard utility
theory inside the formalization; the structural predictions of
Section~\ref{sec:structural_predictions} catalogue what the
formalization makes falsifiable.  The result of this section is the
load-bearing one: the formalization makes the cost of
proxy-instrumental optimization formally bounded, in a way the
qualitative welfare-economics tradition could not.  The shape of
the claim:

\begin{quote}
When the framework optimizes a proxy $P$ for an operational
truth $T$, and the proxy is related to $T$ only instrumentally
($P$ correlates with $T$ but does not constitute it), the
framework's certified alignment-guarantee bound degrades in
proportion to the proxy-to-truth gap $\epsnonres$
(Theorem~\ref{thm:goodhart}).  We additionally conjecture that
under optimization pressure, regions of capability space where
$P$ and $T$ diverge become attractive directions for the
optimizer to exploit, widening the observed $\epsnonres$ at
equilibrium beyond what a less-capable optimizer produces
(Conjecture~\ref{conj:optimization_pressure}, stated and
discussed below; not formally proved in this paper).
\end{quote}

The result is a quantitative version of Goodhart's
Law~\citep{strathern1997improving,manheim2019categorizing} for
capability-space optimization, with the alignment guarantees of the
GFM sequence as its application target.  Two structural choices
shape the result.  First, $T$ is treated as an \emph{operational}
truth measure (the realized exercise contribution
$\volRwin$), not as a metaphysical model of inner preference;
Section~\ref{sec:goodhart_T} makes this explicit.  Second, the
proxy-to-truth gap is split into a within-channel disagreement
component drivable downward by observation, and a residual-class
floor that no amount of within-channel observation can shrink;
Section~\ref{sec:goodhart_gap} states the split.  The theorem itself
(Section~\ref{sec:goodhart_theorem}) is then a Lipschitz-style
continuity result on the within-channel component, with the floor
component contributing a deployment-dependent ceiling on attainable
tightening.

\subsection{The operational truth $T$}
\label{sec:goodhart_T}

The Goodhart machinery requires a truth measure to compare the proxy
against.  We adopt the operational truth measure of
\citep[Revealed~Sacrifice]{lasser2026paper8a}:

\begin{definition}[Operational truth measure]
\label{def:operational_T}
For a capability $c \in P$, the \emph{operational truth measure}
$T(c)$ is the window-active realized exercise contribution
\[
  T(c) \;\equiv\; \volRwin(c).
\]
$T$ is the latent structural quantity that
\citep[Theorem~2]{lasser2026paper8a} bounds from below via the
constructible lower bound $\volRlower$.  Operationally, $T$ is the
quantity an alignment property should track: it is what the
collective \emph{actually} does with capability $c$ during the trade
window $W$, not what the framework's measure says it could do.
\end{definition}

\begin{remark}[Operational, not metaphysical]
\label{rem:T_operational}
$T = \volRwin$ is operational (grounded in the realized-exercise
contribution that revealed-sacrifice trades witness under
S0~(\citep[Revealed~Sacrifice]{lasser2026paper8a},
Section~2.1)), not metaphysical.  The agent's ``true inner
preference,'' considered as an abstract object distinct from
realized exercise, is unobservable in principle, and the framework
does not pretend to access it.  The Goodhart machinery of this
paper operates on operational $T$.  Whether operational $T$
coincides with metaphysical inner preference (in some philosophical
sense the agent might endorse on reflection) is outside the formal
apparatus.
\end{remark}

\begin{remark}[The S0 bridge]
\label{rem:S0_bridge}
The link from the framework's proxy $P = \volL$ to the operational
truth $T = \volRwin$ runs through assumption~S0
of~\citep[Revealed~Sacrifice]{lasser2026paper8a}: an agent's
valuation $\Ui(Y)$ does not exceed the window-active realized
exercise contribution $\volRwin(Y)$.  Without S0, voluntary trade
events do not yield a one-sided constraint relating proxy
disagreement to realized-exercise truth, and the Goodhart machinery
of this paper reverts to a claim about $P$ versus whatever object
anchors $T$, non-vacuous but losing its operational anchor.  S0 is
empirically delicate (typical agents in typical markets do not pay
more than what the capability delivers in exercise) and is the
substance of \citep[Section~2.1]{lasser2026paper8a}.  We take S0 as
standing for the rest of this section.
\end{remark}

\begin{remark}[When operational $T$ is an inadequate truth anchor]
\label{rem:T_failure_modes}
$T = \volRwin$ is operationally adequate as a truth anchor for
the Goodhart machinery in deployments where the trade window
$W$'s recorded sacrifice events are a faithful reflection of
what the agent values, i.e., deployments where the
revealed-preference reading of trades is, in the standard
microeconomic sense, the right interpretation.  Three classes of
deployments where $T = \volRwin$ becomes inadequate as a truth
anchor are worth flagging:
\begin{itemize}
\item \textbf{Asymmetric-information cases.}  When the agent
  trades on private information that systematically diverges
  from the value the trade would carry under symmetric
  information, $\volRwin(Y)$ records the realized exercise
  under the agent's private information state, not under the
  counterfactual ``what would the agent value if their
  information matched the market's?''  Signaling games,
  adverse-selection markets, and information-asymmetric
  insurance markets are canonical examples.  In these
  deployments the operational truth and the
  symmetric-information truth diverge, and the Goodhart bound's
  alignment-relevance depends on which truth the deployment is
  trying to track.
\item \textbf{Time-shifted preferences and identity changes.}
  When the agent's preferences at trade time differ from the
  preferences at exercise time (myopic discounting, addiction,
  identity-shifting interventions), $\volRwin$ records the
  agent's exercise under the post-trade preferences, which may
  not be the preferences the trade was supposed to serve.  S0
  is preserved per-window but the alignment-relevant truth is
  cross-window.
\item \textbf{Coercion and constrained-choice failure of S1.}
  S1 (free choice with non-degrading alternatives) is what
  makes operational $T$ a faithful truth anchor.  When S1
  fails systematically across a deployment (structural
  coercion, market failure, severely constrained choice),
  revealed sacrifice records cost-of-access rather
  than preference, and the operational truth bound on $T$ from
  trade events is correspondingly weakened.
\end{itemize}
The Goodhart machinery of this section presupposes that the
deployment satisfies S0 and S1 sufficiently for $\volRwin$ to
be an adequate operational truth anchor for the
alignment-property bounds the framework is asking about.
Deployments outside that envelope require either a different
truth-anchor convention (with the Lipschitz transfer adjusted
correspondingly) or a separate theory of how the
asymmetry/time-shift/coercion affects the bound.  We close the
loop on what is and is not affected: the Lipschitz bound of
Theorem~\ref{thm:goodhart} \emph{remains valid for the
operational truth as defined}; what degrades under the failure
modes above is the bound's \emph{welfare-relevance}: the
extent to which the operational truth coincides with whatever
counterfactual notion of truth (symmetric-information,
cross-window-stable preference, free-choice baseline) the
deployment ultimately wants to track.  Bounding that secondary
divergence is outside this paper's scope.
\end{remark}

The proxy is unambiguously $P = \volL$: the poset measure the
framework optimizes throughout the
sequence~\citep[Definition~6]{lasser2026gfm}.  The framework
optimizes $\volL$ as if it were the target; the Goodhart machinery
asks what happens when $\volL$ is only instrumentally related to
the operational truth $\volRwin$.

\subsection{The proxy-to-truth gap}
\label{sec:goodhart_gap}

The gap between $P$ and $T$ has two structurally distinct
components, corresponding to two structurally distinct ways the
proxy can fail.  The first is \emph{within-channel disagreement}: on
capabilities the revealed-sacrifice channel can in principle reach,
$P$ and $T$ may differ in magnitude, and that disagreement is what
observation infrastructure can in principle close.  The second is
\emph{channel reach}: capabilities the channel cannot reach
(\citep[Need~Sufficiency]{lasser2026paper8b}, Definition~12,
$\ResS$) contribute a component to alignment slack that no amount
of within-channel observation can shrink.  We define the components
separately and combine them at the end.

\begin{definition}[Within-channel disagreement gap, $\epsnonres$]
\label{def:eps_nonres}
Let $\{c_k\}_{k=1}^K$ enumerate the capabilities in the framework's
poset $P$, and let $\ResS \subseteq P$ be the structural residual
class of \citep[Need~Sufficiency]{lasser2026paper8b}, Definition~12.
The \emph{within-channel disagreement gap} is
\[
  \epsnonres \;\equiv\; \sup_{k \,:\, c_k \notin \ResS,\; T(c_k) > 0}
  \frac{|P(c_k) - T(c_k)|}{T(c_k)},
\]
the worst-case relative deviation of the proxy from the operational
truth on capabilities the sacrifice channel can in principle reach,
restricted to capabilities with strictly positive operational truth.
\end{definition}

\begin{definition}[Residual-class floor, $\epsfloor$]
\label{def:eps_floor}
Capabilities in $\ResS$ are structurally outside the proxy's
measurement scope: trade events do not witness their exercise.
They contribute a \emph{floor} to alignment slack equal to the share
of operational truth the proxy structurally cannot witness:
\[
  \epsfloor \;\equiv\; \frac{\volR(\ResS)}{\volR(P)},
\]
the share of total exercise that lives in the residual class.
$\epsfloor$ is bounded above by the residual-class
characterization of
\citep[Need~Sufficiency]{lasser2026paper8b}, Proposition~8, and
depends on deployment-specific individuation choices about which
capabilities admit ``no sacrifice-observable path in principle.''
\end{definition}

\begin{definition}[Total alignment slack, $\epsgap$]
\label{def:eps_gap}
The \emph{total alignment slack} is
\[
  \epsgap \;\equiv\; \epsnonres + \lambda \cdot \epsfloor,
\]
where $\lambda \in [0, 1]$ weights the residual-class floor's
contribution to alignment slack relative to the within-channel
disagreement.  $\lambda = 0$ corresponds to the framework's ``we
make alignment guarantees only on what the channel can witness''
stance: the residual class is excluded from alignment-quality
accounting.  $\lambda = 1$ corresponds to ``unobservable capability
mass contributes fully to misalignment'' and treats the residual
class as a full liability.  The choice of $\lambda$ is a deployment
question; the main results of this section are stated for
$\epsnonres$, with $\epsfloor$ as a separate floor on attainable
tightening.
\end{definition}

\begin{remark}[Sup-norm conservatism]
\label{rem:supnorm_conservative}
The sup-norm in Definition~\ref{def:eps_nonres} means a single
badly-proxied dimension controls the disagreement gap.  This is
conservative but correct for alignment purposes: one dimension
where the proxy is misleading is sufficient to produce a misaligned
decision when the framework optimizes the proxy.  A weighted norm
that gives less weight to safety-irrelevant dimensions could give
a less conservative measure, but the safety-relevance weighting is
itself application-dependent and not addressed in this paper; the
sup-norm is the worst-case-but-defensible choice.
\end{remark}

\begin{remark}[Aggregate diagnostic ceiling vs sup-norm bound]
\label{rem:agg_vs_sup}
$\epsnonres$ is defined as a sup-norm.  Operationally, the
framework does not in general have access to per-dimension
estimates of $|P(c_k) - T(c_k)|/T(c_k)$; what it has access to via
\citep[Theorem~2]{lasser2026paper8a} is the aggregate diagnostic
ceiling $(1 - \betaL)/\betaL$ derived from the aggregate lower
ratio $\betaL$, computed at the bundle or per-capability
level depending on Part~A vs Part~B aggregation.  The aggregate
ratio bounds an \emph{aggregate} version of $\epsnonres$ from
above; it bounds the sup-norm $\epsnonres$ only if per-capability
lower-ratio assumptions hold.  We use $(1 - \betaL)/\betaL$ as a
\emph{tractable diagnostic ceiling} on attainable tightening
under aggregate observation, not as a sup-norm bound.  When the
deployment supports per-capability observation (e.g., bundle
decomposition, hedonic-disaggregation panels), $\epsnonres$ can be
bounded directly per-dimension.
\end{remark}

\begin{remark}[Cell-wise attribution of the diagnostic ceiling]
\label{rem:cellwise_attribution}
The five-cell $\delta$-decomposition of
\citep[Need~Sufficiency]{lasser2026paper8b}, Definition~8,
attributes the diagnostic ceiling across cells:
$\delta_{\mathrm{covered}}$ is the within-channel
proxy-failure-candidate component (contributing to $\epsnonres$);
$\delta_{\mathrm{dormant}}$, $\delta_{\mathrm{restricted}}$, and
$\delta_{\mathrm{boundary}}$ are structural-invisibility components
on capabilities the channel could reach but has not yet (these
contribute to $\epsnonres$ in proportion to their pending
observation share); $\delta_{\mathrm{residual}}$ is the
$\epsfloor$ component proper.  The cell decomposition is what the
deployment uses to identify \emph{which} kind of intervention
(more trades, sharper attestation, finer individuation) is the
operative tightener for the deployment's actual gap.
\end{remark}

\begin{remark}[Floor non-zero, but not unique to GFM]
\label{rem:floor_shared}
$\epsfloor$ is non-zero in any deployment with non-empty $\ResS$.
The fungibility-collapse limit of
Section~\ref{sec:fungibility_collapse} drives $\epsnonres$ on the
observed-and-non-residual subset toward zero (the proxy is exact
when the poset is fungible), but does \emph{not} drive $\epsfloor$
to zero.  Standard economics has the same floor, expressed
informally as the shadow-economy / non-market-activity gap that
imputation methods (imputed rent, time-use surveys, satellite
accounts) only partially close~\citep{stiglitz2009report}.
$\epsfloor$ is the GFM-side formal counterpart to this shared
problem; the discussion in Section~\ref{sec:lineage} (the SSF
subsection) traces the lineage in detail.
\end{remark}

\subsection{The Goodhart theorem}
\label{sec:goodhart_theorem}

\begin{theorem}[Goodhart Lipschitz transfer]
\label{thm:goodhart}
Define the \emph{operationally active subspace}
\[
  P^{\mathrm{act}} \;\equiv\; \{c \in P \setminus \ResS \,:\, T(c) > 0\},
\]
the set of capabilities outside the residual class that the
operational truth witnesses as exercised in the trade window~$W$.
Let $g(\cdot)$ be an alignment property in the GFM sequence,
viewed as a real-valued function of a capability measure
restricted to~$P^{\mathrm{act}}$.  If $g$ is Lipschitz on this
restriction with property-specific Lipschitz constant
$\mathrm{Lip}(g)$ relative to the sup-norm of relative deviation,
then under $\volL$-maximization with proxy-to-truth disagreement
gap $\epsnonres > 0$,
\[
  \bigl| g(T) - g(P) \bigr| \;\leq\; \mathrm{Lip}(g) \cdot \epsnonres.
\]
That is, the alignment property's value under the operational
truth $T$ deviates from the value under the proxy $P$ by at most
$\mathrm{Lip}(g) \cdot \epsnonres$.
\end{theorem}

\begin{conjecture}[Optimization pressure and gap exploitation]
\label{conj:optimization_pressure}
Let an optimizer with capability measure $C$ act on the proxy
$P$.  Regions of capability space where $P$ and $T$ diverge are
\emph{available} directions for the optimizer to move into;
cheaper $P$-gain per unit $T$-loss makes these regions attractive
relative to regions where $P$ and $T$ agree.  We conjecture that:
\begin{quote}
More capable optimizers, with greater reach across capability
space, have access to a larger set of $P$-$T$-divergence regions;
an optimizer that exploits any of them widens the observed
$\epsnonres$ at equilibrium beyond what a less-reaching optimizer
produces.
\end{quote}
The conjecture is qualitative.  No formal proof is given in this
paper; a precise functional form $f: C \mapsto \epsnonres(C)$ is
optimizer-architecture-dependent and is the subject of
Open~Question~\ref{oq:goodhart_rate}.
\end{conjecture}

\begin{remark}[Theorem 2 is conditional; Conjecture 1 is unproven]
\label{rem:thm2_qualitative}
Theorem~\ref{thm:goodhart} is conditional on $g$ being Lipschitz
on $P^{\mathrm{act}}$ with a specific constant $\mathrm{Lip}(g)$;
the per-property Lipschitz analyses (computing
$\mathrm{Lip}(g)$ for the anti-monopolar threshold, the phase
boundary, the convergence rate, the damage bound, and the
aggregate-lower-bound slack) are
\emph{perturbative roadmap items}, not theorems of the present
paper.  Conjecture~\ref{conj:optimization_pressure} is a
directional claim about the relationship between optimizer
capability and observed $\epsnonres$ at equilibrium; we do not
prove it here.  The present paper establishes the conditional
Lipschitz transfer and identifies the optimization-pressure
direction as conjectural pending a formal optimizer model
(Open~Question~\ref{oq:goodhart_rate}).
\end{remark}

\begin{proof}[Proof of Theorem~\ref{thm:goodhart}]
The argument is a Lipschitz-continuity transfer.  Let
$g: \mathcal{M} \to \mathbb{R}$ be the alignment property, where
$\mathcal{M}$ is the space of capability measures restricted to
the operationally active subspace $P^{\mathrm{act}}$.  By
hypothesis, $g$ is Lipschitz with constant $\mathrm{Lip}(g)$ in
the sup-norm of relative deviation: for any two measures
$\mu, \nu \in \mathcal{M}$ with $\nu(c_k) > 0$ for all $c_k \in
P^{\mathrm{act}}$,
\[
  \bigl| g(\mu) - g(\nu) \bigr| \;\leq\;
  \mathrm{Lip}(g) \cdot \sup_{c_k \in P^{\mathrm{act}}}
  \frac{|\mu(c_k) - \nu(c_k)|}{\nu(c_k)}.
\]
Setting $\mu = P$ and $\nu = T$ (with $T(c_k) > 0$ on
$P^{\mathrm{act}}$ by construction), the right-hand sup-norm
equals $\epsnonres$ exactly:
\[
  \sup_{c_k \in P^{\mathrm{act}}}
  \frac{|P(c_k) - T(c_k)|}{T(c_k)} \;=\; \epsnonres.
\]
Combining,
\[
  \bigl| g(T) - g(P) \bigr| \;\leq\; \mathrm{Lip}(g) \cdot \epsnonres,
\]
which is the claim.
\end{proof}

\begin{remark}[Capabilities with $T(c) = 0$ are evidence-inactive in $W$]
\label{rem:zero_T_inactive}
The restriction of $g$ to $P^{\mathrm{act}}$ in the theorem
statement is structurally important.  A capability $c \in
P \setminus \ResS$ with $T(c) = 0$ is channel-reachable but
\emph{operationally inactive} for the alignment guarantees of
this paper: the framework knows it exists (it lies outside
$\ResS$) but the operational truth $T = \volRwin$ records no
exercise contribution from it within~$W$.  Capabilities in
$\delta_{\mathrm{dormant}}$ or $\delta_{\mathrm{restricted}}$
of \citep[Definition~8]{lasser2026paper8b} are
\emph{evidence-inactive} in $W$ (no admissible sacrifice events
witness their exercise) but may carry latent
$\volR^{[W]}$ that the channel has not yet attested; whether
evidence-inactive capabilities are also operationally inactive
($T(c) = 0$ in the strict sense the bound takes) is a deployment
accounting choice.  A strict reading sets $T(c) = 0$ on
evidence-inactive capabilities, while a relaxed reading admits
imputed $T(c)$ from non-trade evidence at the cost of leaving
the Goodhart bound's S0 anchor.  Channel~1 (observation density)
is the operational lever that brings evidence-inactive
capabilities into $P^{\mathrm{act}}$ over time, by accumulating
admissible events that record their exercise.  Until that
happens, the Goodhart bound's content on these capabilities is
vacuous under the strict reading: any deviation
$|P(c) - T(c)|$ on a capability with $T(c) = 0$ is, in the trade
window~$W$, an operationally inactive deviation.  The Lipschitz
hypothesis on $P^{\mathrm{act}}$ is the right scope for the
bound to be non-trivial; on the complement, the bound is
trivially zero on the truth side and Channel~1 brings new
capabilities under the
bound's content as observation accumulates.
\end{remark}

\begin{remark}[Structural argument supporting Conjecture~\ref{conj:optimization_pressure}]
\label{rem:conj1_structural}
The structural argument supporting
Conjecture~\ref{conj:optimization_pressure}, sketched here, is
not a proof.  An optimizer that maximizes $P$ subject to
dynamics on capability space will, at equilibrium, satisfy
$\partial P / \partial \mathrm{action} = 0$, but
$\partial T / \partial \mathrm{action}$ need not be zero where
$P$ and $T$ diverge.  This creates a ``free direction'' the
optimizer can exploit.  The size of this free direction grows
with the optimizer's reach in capability space, in the sense
that more capable optimizers have access to a larger set of
capability-space regions where the divergence is exploitable.
A formal version of the argument requires a model of how the
optimizer searches capability space (greedy local,
gradient-based, planning-based) and is correspondingly
architecture-dependent.  Open Question~\ref{oq:goodhart_rate}
of Section~\ref{sec:discussion} states the precise question.
\end{remark}

\begin{remark}[Empirical witness, not causal proof]
\label{rem:hhi_empirical_witness}
The wireheading-consistent HHI of
\citep[Need~Sufficiency]{lasser2026paper8b}, Propositions~6 and 7,
provides a deployment-readable empirical witness for the
mechanism of Conjecture~\ref{conj:optimization_pressure}: trade-flow
concentration in a narrow set of bundle categories is a signature
\emph{consistent with} an optimizer exploiting $P$-$T$
divergence.  But the witness is correlational, not
causal-inferential.  Other mechanisms also produce concentration:
market consolidation, demand shocks, supply-side limitations.
Distinguishing the optimization-pressure mechanism from these
alternatives requires natural-experiment or instrumental-variable
designs that are outside the framework's structural apparatus.
The HHI is a deployment-readable indicator that warrants further
investigation when crossed; it does not by itself establish that
optimization pressure is the cause.  As
\citep[Propositions~6~and~7]{lasser2026paper8b} establish for the
HHI itself, ``wireheading-consistent'' is the right operational
phrasing.
\end{remark}

\subsection{Candidate alignment properties: the perturbative roadmap}
\label{sec:goodhart_examples}

The conditional Lipschitz-transfer of Theorem~\ref{thm:goodhart}
applies to any alignment property $g$ that admits a Lipschitz
analysis on the relevant capability subspace.  Five such properties
are natural candidates from the GFM sequence; we list them as
\emph{perturbative roadmap items} (candidate targets for a
subsequent perturbation analysis), not as bounds the present
paper establishes.

\begin{itemize}
\item \textbf{Anti-monopolar threshold.}  The threshold
  $\gamma^*$ at which the anti-monopolar correction flips sign
  (\citep[Proposition~6]{lasser2026horizon}) is a function of the
  capability measure on a multi-substrate poset.  The expected
  perturbative bound:
  $|\gamma^*_{T} - \gamma^*_{P}| \leq C_{\gamma} \cdot \epsnonres$
  for some $C_{\gamma}$ derived from the property's structural
  sensitivity to substrate-volume distributions.  Establishing
  $C_{\gamma}$ is a perturbation-analysis question on a specific
  family of multi-substrate measures.

\item \textbf{Phase-boundary location.}  The self-correcting basin
  of \citep[Theorem~1]{lasser2026phase} is a region of joint state
  space defined by Lyapunov-function-style stability conditions on
  $\volL$.  The expected perturbative result: under the truth $T$,
  the basin differs from the proxy-defined basin by a boundary
  layer of width $O(\epsnonres)$, with the constant depending on
  the Lyapunov function's smoothness.

\item \textbf{Convergence rate.}  The KL-rate of
  \citep[Theorem~2]{lasser2026wamura} is a function of the
  proxy-truth discrepancy in the relaxation-protocol's
  inferential bookkeeping.  The expected perturbative result: the
  convergence rate under $T$ deviates from the rate under $P$ by
  $O(\epsnonres)$, with sign deployment-dependent (the proxy may
  predict either faster or slower convergence than the truth).

\item \textbf{Damage bound.}  The damage bound $\damagebound$ of
  \citep[Theorem~1]{lasser2026wamura} bounds the contraction of
  $\volL$ under controlled relaxation.  Under the truth $T$, the
  damage to operational exercise differs from the proxy-bounded
  damage by $O(\epsnonres \cdot \max\text{-contraction}(\scope))$,
  where $\scope$ is the controlled-relaxation scope.

\item \textbf{Aggregate-lower-bound slack.}  $\betaL$ from
  \citep[Theorem~2]{lasser2026paper8a} is the aggregate
  \emph{lower ratio} of the construction; the corresponding
  diagnostic ceiling on $\epsnonres$ is $(1 - \betaL)/\betaL$.
  As $\epsnonres$ shrinks, $\betaL$ approaches the true
  $\beta = T_{\mathrm{observed}} / T_{\mathrm{total}}$ at a rate
  that decomposes into S0 attestation slack, Part-A vs Part-B
  aggregation, and observation density.
  This is the most direct connection between Theorem~\ref{thm:goodhart}
  and the measurement theorem of
  \citep[Revealed~Sacrifice]{lasser2026paper8a}: the framework's
  operational tightening tool drives down the same gap the Goodhart
  theorem identifies as alignment-controlling.
\end{itemize}

The per-property Lipschitz constants are not established by this
paper's machinery and are listed only to indicate where the
conditional transfer result lands.  Subsequent perturbation
analysis closes each per-property bound under property-specific
sensitivity assumptions.

\subsection{Better economics tightens alignment: the corollary}
\label{sec:goodhart_corollary}

The conditional Lipschitz-transfer of Theorem~\ref{thm:goodhart}
combined with the monotonicity of the revealed-sacrifice channel
gives the paper's headline corollary.

\begin{corollary}[Direction: better economics tightens alignment]
\label{cor:direction}
Suppose the revealed-sacrifice channel of
\citep[Revealed~Sacrifice]{lasser2026paper8a} is admissibly
accumulating observations, and assume per-capability
admissibility (the deployment satisfies the conditions formalized
in Assumption~\ref{ass:per_cap_admissibility} of
Section~\ref{sec:lineage}, which bridges the aggregate diagnostic
ceiling to the sup-norm $\epsnonres$).  Then:
\begin{enumerate}
\item By
  \citep[Propositions~2~and~3]{lasser2026paper8a}, $\betaL$ is
  non-decreasing in the observation count, so the
  diagnostic ceiling $(1 - \betaL)/\betaL$ on the aggregate
  version of $\epsnonres$ is non-increasing.
\item Under per-capability admissibility, the diagnostic
  ceiling bounds the sup-norm $\epsnonres$, and by
  Theorem~\ref{thm:goodhart} the \emph{certified upper bound}
  on each Lipschitz alignment property's deviation from the
  operational truth is non-increasing in the same direction.
  (The actual deviation $|g(T) - g(P)|$ is what it is at any
  given moment; what tightens monotonically is the bound the
  framework can certify.)
\item By Definition~\ref{def:eps_gap}, the total alignment slack
  $\epsgap = \epsnonres + \lambda \cdot \epsfloor$ is bounded
  below by $\lambda \cdot \epsfloor$, so tightening saturates at
  the deployment-specific residual-class floor.
\end{enumerate}
That is: improving the economic model (more trade data, better
bundle disaggregation, wider coverage, sharper S0 attestation,
finer five-cell $\delta$-attribution) tightens each Lipschitz
alignment guarantee in the direction of the operational truth, with
saturation at the deployment-specific $\lambda \cdot \epsfloor$
floor.
\end{corollary}

\begin{remark}[Direction, not rate]
\label{rem:direction_not_rate}
Corollary~\ref{cor:direction} states a \emph{direction}
(monotonicity), not a \emph{rate}.
\citep[Propositions~2~and~3]{lasser2026paper8a} are existence
statements about non-decreasing $\betaL$; they do not bound how
fast $\betaL$ converges, and the rate depends on
(i) S0 attestation quality (which itself is institutional-layer-
dependent;~\citep[Remark~18 (external attestation)]{lasser2026paper8a}),
(ii) Part-A vs Part-B aggregation choice
(\citep[Theorem~2 Part~A and Part~B]{lasser2026paper8a}),
(iii) observation density on the deployment's capability
inventory, and (iv) the residual-class floor $\epsfloor$ itself.
None of (i)--(iv) is bounded in this paper.  Open
Question~6 of Section~\ref{sec:discussion} states the
quantitative-rate question precisely; until it is bounded, the
corollary is restricted to direction.
\end{remark}

Corollary~\ref{cor:direction} states a one-mechanism direction.  The
strong form of the better-economics-tightens-alignment claim,
together with the channel-wise decomposition that gives the
direction operational substance, lives in
Section~\ref{sec:lineage}.

\subsection{Position relative to standard Goodhart literature}
\label{sec:goodhart_lineage}

The result of this section is recognizable as a structural variant
of Goodhart's Law~\citep{strathern1997improving} for the
specific case of capability-space optimization in the GFM
framework.  The categorization of \citep{manheim2019categorizing}
distinguishes four Goodhart variants: regressional, extremal,
causal, and adversarial.  The mapping from this paper to that
classification is clean.  Theorem~\ref{thm:goodhart} is closest
to the \emph{regressional} variant: the proxy and truth
correlate but are not identical, and the worst-case
active-subspace Lipschitz transfer bounds how much correlation
imperfection propagates to alignment-property deviations under
proxy-instrumental optimization.
Conjecture~\ref{conj:optimization_pressure} is closest to the
\emph{adversarial} variant: a sufficiently capable optimizer
treats $P$-$T$ divergence regions as exploitation targets,
widening the observed gap at equilibrium.  The other two
variants (\emph{extremal}, where the proxy-truth relation
breaks at distribution tails, and \emph{causal}, where
optimizing the proxy disrupts the causal mechanism that made
it a proxy in the first place) are not directly addressed by
this paper's results, though both arise naturally in
deployment settings the framework would extend to.

Recent formal work on Goodhart's law in AI alignment frames the
phenomenon through tail-distribution
analyses~\citep{elmhamdi2024goodhart,majka2025strong},
complementing the worst-case active-subspace Lipschitz framing
of Theorem~\ref{thm:goodhart}: the tail-distribution treatments
characterize how proxy-truth deviation behaves on extreme events
in distribution, while the Lipschitz transfer here gives a
worst-case bound across the active-subspace.  The two
formalizations are not in tension; they describe complementary
slices of the same underlying mechanism (extremal-event behavior
versus active-subspace worst case), and a complete theory of
proxy-instrumental alignment would compose both.

The novel contribution beyond the literature's framing is the
structural embedding: in GFM, $P = \volL$ and $T = \volRwin$
are not abstract ``proxy'' and ``truth''; they are formally
specified objects with axiomatic structure, and the proxy-truth
gap admits the within-channel + residual-class decomposition of
Section~\ref{sec:goodhart_gap}.  This decomposition is what lets
the corollary bind tightening to specific operational
interventions (more trades, sharper attestation, finer
individuation) rather than to abstract ``proxy quality
improvement.''
